% sections/05_aggregation.tex
\section{Aggregation and Classification}\label{sec:aggregation}

\subsection{Aggregation Pipeline}\label{sec:agg:pipeline}

The end-to-end computation from raw data to published output follows six
sequential stages:

\begin{enumerate}
  \item \textbf{Share-vector construction.}  For each country--axis pair (and,
    where applicable, each channel or fuel category), partner-level flows are
    normalised to sum-to-one share vectors.
  \item \textbf{Concentration scoring.}  The concentration formula
    (\cref{eq:hhi}) is applied to each share vector, yielding channel-level,
    fuel-level, or mode-level scores as appropriate for the axis (see
    \cref{sec:axis-specifications}).
  \item \textbf{Axis-level aggregation.}  Channel scores are combined via
    the arithmetic mean (\cref{eq:dual-channel}), fuel scores via the fuel
    average (\cref{eq:fuel-avg}), or tonnage-weighted scores via the
    category-weighted formula (\cref{eq:weighted-b}), according to the
    axis-specific variant.
  \item \textbf{Composite averaging.}  The six axis scores are
    arithmetically averaged to produce the composite~\(\mathrm{ISI}_i\)
    (\cref{eq:composite}).
  \item \textbf{Rounding.}  All published scores---axis-level and
    composite---are rounded to eight decimal places
    (\texttt{ROUND\_PRECISION\,=\,8}, round-half-to-even).  Rounding is
    applied before classification (step~6) and before hashing
    (\cref{sec:integrity:country}).  Classification operates exclusively
    on rounded values.
  \item \textbf{Classification.}  The rounded composite is mapped to one of
    four ordinal tiers (\cref{tab:thresholds}) by a single classification
    function in the classification module.
\end{enumerate}

No step in this pipeline involves discretionary judgement.  Given identical
input data, the output is fully deterministic.

\subsection{Missing-Axis Protocol}\label{sec:agg:missing}

If axis~\(a\) is non-computable for country~\(i\) (due to null channels,
zero flows, or data-coverage failure), the composite is computed over the
remaining axes:
%
\begin{equation}\label{eq:composite-partial}
  \mathrm{ISI}_i
  \;=\;
  \frac{1}{|\mathcal{A}_i|}
  \sum_{a \in \mathcal{A}_i} M_i^{(a)}
\end{equation}
%
where \(\mathcal{A}_i \subseteq \{1, \ldots, 6\}\) is the set of computable
axes for country~\(i\).  If \(|\mathcal{A}_i| < 4\), the country is excluded
from the ranked output and flagged \texttt{W-COV}.

\paragraph{v\,1.0 status.}
In the \ISI v\,1.0 snapshot (EU-27, vintage~2024), all six axes are computable
for all 27~countries.  The denominator in \cref{eq:composite} is therefore
fixed at~6 for every country in the published output.  The fallback rule
defined above is not triggered in v\,1.0 but remains part of the versioned
methodology specification.

\subsection{Ranking Methodology}\label{sec:agg:ranking}

Countries are ranked in descending order of the composite score.  Ties are
broken lexicographically by ISO~3166-1 alpha-2 code.  The ranking is ordinal;
no cardinality is attached to rank differences.

\subsection{Vintage and Periodicity}\label{sec:agg:vintage}

Each published snapshot carries a vintage year~\(t\).  The vintage refers to
the most recent data year for which the majority of axes have coverage.  Due
to source-specific publication lags (see \cref{sec:data:constraints}), a
snapshot published in calendar year~\(t + 1\) or~\(t + 2\) may carry vintage
year~\(t\).  For v\,1.0, the vintage year is \textbf{2024}.

The vintage year is recorded in the snapshot metadata and in the integrity
chain (\cref{sec:integrity-chain}).